\section{Recovering the cut on exceptional episodes}
\label{decode:section}

Every episode in $D$ satisfies $Z(t)\lor A_*(t,q)$. This is useful
information for a suffix: if its prefix ended at $k=s+p$, the true origin
is one of the finitely many candidates $k+L-p'$, $p'\in\mathcal P$.
We will use the exceptional predicate to identify it.

Two issues must be resolved separately. First, false candidate origins
must be sparse over \emph{all} stencil slots, not just those whose
metadata might be correct. Second, a suffix proposing an incorrect end
must not make the true witness disappear. We first count using actual
episode data. We then give a suffix test that obtains precisely enough
of those data before it asks for uniqueness, and checks the decoded
origin's end afterward.

\subsection{A bounded list of possible clock anchors}

Fix a complete episode $[s,b)$ and its actual origin $t=s+L$.
For every pair $p,p'\in\mathcal P$, set
\[
       \delta=p-p',\qquad \tau=t+\delta.
\]
The distinct-slot differences are distinct as integers and modulo
$B$, are separated from one another by more than $2w_0+2$, and
satisfy $|\delta|\le h$.  The identical-slot pairs all give $t$.
Evaluate the bounded larger end search $\zeta(\tau)$ on the
actual episode.

If one of these larger searches is absent, then every smaller
search $\eta(\tau)$ in this collection is absent.  Indeed, a
marker eligible for a smaller search at any one origin has
distance at most $\max_u D^0_u+2h$ from the end-search center at
any other origin.  The inequality
\begin{equation}
                 \max_u D^0_u+2h<\min_u D^1_u
 \label{decode:two-searches}
\end{equation}
makes it eligible for that larger search, whatever residue index
it has at the latter origin.  This would contradict the absent
larger search.

In this absence case the seeds at all these origins have short
periods.  Their starting positions differ by at most $2h$; the
seed length and marker separation ensure an overlap of at least
$d_\bullet^2$.  Lemma~\ref{marker:periodic} gives the same
two-sided periodic continuation for all seeds.  They agree with
the episode on their union, so their first subsequent mismatch
is the same position $q_0$.  This is the actual episode's mismatch
anchor, because its own smaller search is absent.  Define
\begin{equation}
 Q_\tau=
 \begin{cases}
  \zeta(\tau),&\zeta(\tau)\ne\bot,\\
  q_0,&\zeta(\tau)=\bot.
 \end{cases}
 \label{decode:anchors}
\end{equation}
The second case is used only when a larger search is absent,
so the preceding argument always supplies its $q_0$.
In particular $Q_t=q$.

\begin{lemma}[Number of possible anchors]\label{decode:anchor-list}
For a fixed actual episode, the values $Q_\tau$ in
\eqref{decode:anchors}, over all $p,p'\in\mathcal P$, belong to a
set $\mathcal Q$ of at most $20$ positions.
\end{lemma}

\begin{proof}
Every present $\zeta(\tau)$ is an end marker in
\[
 [z_\bullet(t)-h-\max_uD^1_u,
   z_\bullet(t)+h+\max_uD^1_u].
\]
Since $\max_uD^1_u=3d_\bullet+(B-1)S_\bullet$ and
$d_\bullet>10BS_\bullet+10(h+w+1)$, this interval has length
strictly less than $(31/5)d_\bullet$.  End markers have separation
at least $d_\bullet$, so it contains at most seven of them.
There is at most one additional value, the shared mismatch $q_0$.
In particular the asserted bound of $20$ holds.
\end{proof}

\subsection{Counting every false origin in the stencil}

We now count the ordered pairs with $p\ne p'$ for which the
translated origin satisfies the exceptional predicate.  All
marker decisions can be made in one common two-sided extension:
the margin construction makes every inspection used below
internal to the episode, including the shifted searches.

\begin{proposition}[Whole-stencil sparsity]\label{decode:sparsity}
For each complete episode and the positions $Q_\tau$ in
\eqref{decode:anchors},
\begin{equation}
 \#\bigl\{(p,p')\in\mathcal P^2:p\ne p',\
 Z(t+p-p')\ \text{or}\
 A_*(t+p-p',Q_{t+p-p'})\bigr\}<\mu.
 \label{decode:whole-stencil-bound}
\end{equation}
This bound includes every ordered distinct-slot pair, independently
of the kinds assigned to its slots.
\end{proposition}

\begin{proof}
We separately count end-search instability, and the three causes
in $A_*$.  Write $\delta=p-p'$ and $\tau=t+\delta$.

\paragraph{End-search instability costs at most four pairs.}
If an end search changes when $\tau$ is replaced by $\tau+v$,
where $v\in B\mathbb Z$ and $|v|\le w$, the eligibility of some
marker changes.  Otherwise the set of eligible markers, and
hence its leftmost member or its absence, would be unchanged.
For that marker $m_e$, the index
\[
                 u=(m_e-z_\bullet(t)-\delta)\bmod B
\]
is unchanged during this shift, since $v$ is divisible by $B$.
Thus for some $\sigma\in\{0,1\}$ and one of the two signs,
\begin{equation}
 \bigl|m_e-(z_\bullet(t)+\delta\mathbin{\pm}D^\sigma_u)\bigr|
 \le w.
 \label{decode:end-crossing}
\end{equation}
Fix the search $\sigma$ and the sign.  As $\delta$ varies in
$[-h,h]$ and $u$ varies in $\{0,\ldots,B-1\}$, all possible
markers in \eqref{decode:end-crossing} lie in an interval of
diameter at most
\[
                 2h+(B-1)S_\bullet+2w<d_\bullet.
\]
There is at most one end marker in that interval.  For this fixed
marker the indices $u$ are distinct for the distinct-slot
differences, because those differences are distinct modulo $B$.
Two target positions in \eqref{decode:end-crossing} are then
separated by at least $S_\bullet-2h>2w$.
At most one pair can satisfy \eqref{decode:end-crossing} for this
search and sign.  The two searches and two signs therefore
contribute at most four pairs in total.

\paragraph{Fixing an anchor makes the clock count uniform.}
For each fixed $Q\in\mathcal Q$, the clock part of $A_*(\tau,Q)$
is
\[
 \varphi(Q-t)+\varphi(p')-\varphi(p)\in W.
\]
Apply Lemma~\ref{clock:robust} at the translate
$\varphi(Q-t)$, with the realized roster values $\varphi(p)$.
It counts at most $c_{\rm cl}\log(2m)$ ordered distinct-slot
pairs.  Its uniformity in the translate permits this application
for every actual value of $Q$, including an unbounded mismatch
position.  There are only $|\mathcal Q|\le20$ such applications.

\paragraph{For each fixed anchor and label, inner-marker failures
cost at most three pairs.}
Fix $Q\in\mathcal Q$ and $j\in\{0,\ldots,g\}$.
For each $\delta$ write
\[
 z=z_j(t)+\delta,\qquad u=(Q-t-\delta)\bmod B.
\]
In the stability part of $A_*$, the index $u$ is held fixed while
$z$ is shifted by an integer of magnitude at most $w_0$.
The union of all these possible left-endpoint positions has
diameter $2h+2w_0<d$, so it contains at most one inner marker.
The distinct differences $\delta$ are separated by more than
$2w_0+2$.  Consequently at most one pair has its left endpoint
within $w_0$ of this marker.  Charge that pair, if present, to
the left endpoint.

For any remaining pair, no marker crosses the left endpoint
during the allowed shifts.  The first marker on or after $z$,
if it exists, is therefore fixed during those shifts.
A change of the first-marker search can now occur only if this
marker is within $w_0$ of the right endpoint $z+H_u$.
Also, if the search on $[z,z+H_u]$ is absent but the extended
interval contains a marker, its first marker lies in
$(z+H_u,z+H_u+H_{\rm ext}]$.
Thus either remaining failure puts the first marker on or after
$z$ within $\max(H_{\rm ext},w_0)$ of $z+H_u$.

As $z$ ranges over its sweep of diameter $2h<d$, at most two
positions can play this first-marker role: a marker within the
sweep, and the first marker after the sweep.  If no such marker
exists, there is no pair to charge to it.
For distinct-slot differences the indices
$u=(Q-t-\delta)\bmod B$ are all distinct.  The spacing of the
widths in the schedule gives, for different such differences,
\begin{align*}
 |(z_j(t)+\delta+H_u)-(z_j(t)+\delta'+H_{u'})|
 &\ge |H_u-H_{u'}|-|\delta-\delta'|\\
 &>2\max(H_{\rm ext},w_0)+2.
\end{align*}
Each of the two possible first markers can therefore account for
at most one pair.  Together with the possible left-endpoint pair,
this is at most three pairs for fixed $Q,j$.

Combining the counts, the left side of
\eqref{decode:whole-stencil-bound} is at most
\begin{equation}
 4+20\bigl(c_{\rm cl}\log(2m)+3(g+1)\bigr)
 <20\bigl(c_{\rm cl}\log(2m)+4(g+1)\bigr)+10<\mu.
 \label{decode:budget}
\end{equation}
The last inequality is precisely the multiplicity choice in
the schedule.  Counting all pairs for each fixed $Q$, rather
than only pairs assigned that $Q$, is an upper bound and makes
the use of the clock and width estimates uniform.  Taking the
union over the short list $\mathcal Q$ completes the proof.
\end{proof}

There are two levels to this estimate. At a \emph{fixed} anchor $Q$,
distinct displacement residues give distinct width indices, so right-end
failures are sparse. The anchor list then permits at most twenty such
counts. The estimate includes all ordered pairs; no knowledge of the
correct kind was used to obtain it. We can now turn it into a decoder
without assuming that a suffix's proposed endpoint is correct.

\subsection{A decoder that also verifies its endpoint}

We next give an actual test on a suffix beginning at $k$.
For each slot $p'\in\mathcal P$ form the candidate origin
\begin{equation}
                        \tau_{p'}=k+L-p'.
 \label{decode:candidate-origins}
\end{equation}
These are positions relative to the suffix; the origin itself
need not be a visible letter.  All of the following inspections,
including every shifted search needed to evaluate $Z$ and $A_*$,
are mandatory.  If a required window is absent from the suffix,
the test rejects.

First compute every $\zeta(\tau_{p'})$.  If at least one is
absent, require all $\eta(\tau_{p'})$ to be absent, read every
candidate seed, and require the suffix to end exactly according
to the end rule at one fixed reference candidate.  This includes
agreement with that reference seed's periodic continuation up
to its first subsequent mismatch, and exactly $L_{\rm tail}$
letters from that mismatch to the proposed endpoint.  Denote
its verified mismatch position by $q_{\rm ref}$.  For candidate
origins with absent $\zeta$, use $Q_{\tau_{p'}}=q_{\rm ref}$.
For every candidate with present $\zeta$, always use
$Q_{\tau_{p'}}=\zeta(\tau_{p'})$.
If all the larger searches are present, no reference mismatch
test is needed at this step; all these $Q$'s are bounded search
returns.

Now require that exactly one slot $p'$ satisfy
\begin{equation}
         Z(\tau_{p'})\ \text{or}\
         A_*(\tau_{p'},Q_{\tau_{p'}}).
 \label{decode:unique-witness}
\end{equation}
This is uniqueness over all slots, without any restriction to
particular kinds or guessed table values.  Finally, after
decoding this slot, require exact consumption according to the
end rule at its candidate origin in \emph{all} cases, including
the case where all larger searches were present.

Figure~\ref{decode:two-guards-figure} separates the two anchor roles and
records the order of these end checks. The reference guard belongs only
to the branch with an absent larger search. The decoded origin's exact
end is checked in both branches.
\begin{figure}[htbp]
\centering
\begin{tikzpicture}[
  font=\small,
  box/.style={draw,rounded corners=2pt,align=center,inner sep=5pt},
  flow/.style={-{Stealth[length=2mm]},semithick},
  every node/.style={outer sep=0pt}
]
\node[anchor=west,font=\small\bfseries] at (-6.45,0)
  {When the end rule succeeds at one origin};
\node[box,text width=6.05cm,minimum height=2.2cm] (endrole)
  at (-3.25,-1.45)
  {\textbf{Smaller search $\eta$: episode end}\\[2pt]
   $q_0=\eta$ if present; otherwise $q_0$ is\\
   the first mismatch after the seed.\\[2pt]
   The exact endpoint is $b=q_0+L_{\rm tail}$.};
\node[box,text width=6.05cm,minimum height=2.2cm] (clockrole)
  at (3.25,-1.45)
  {\textbf{Larger search $\zeta$: clock anchor}\\[2pt]
   $q=\zeta$ if present; otherwise $q=q_0$.\\[2pt]
   A present $\zeta$ need not equal $\eta$.};
\begin{scope}[yshift=-3mm]
\node[anchor=west,font=\small\bfseries] at (-6.45,-2.85)
  {Suffix decoder: candidates $\tau_{p'}=k+L-p'$, for every $p'\in\mathcal P$};
\node[box,text width=11.7cm] (read) at (0,-3.58)
  {Read every bounded search and inspection required by the decoder.\\
   Missing required data cause rejection. Compute all $\zeta(\tau_{p'})$.};
\node[box,text width=5.9cm,minimum height=3.5cm] (present)
  at (-3.25,-6.35)
  {\textbf{Every larger search is present}\\[4pt]
   Use $Q_\tau=\zeta(\tau)$ for every candidate.\\[4pt]
   These clock anchors are bounded.\\
   No reference mismatch test is needed.};
\node[box,text width=5.9cm,minimum height=3.5cm] (absent)
  at (3.25,-6.35)
  {\textbf{Some larger search is absent}\\[4pt]
   Require every $\eta(\tau)=\bot$ and all seeds.\\
   Check one fixed reference origin's exact end\\
   and its first mismatch $q_{\rm ref}$.\\[2pt]
   Use $Q_\tau=\zeta(\tau)$ if present;\\
   otherwise use $Q_\tau=q_{\rm ref}$.};
\draw[flow] (read.south) -- ++(0,-.23) -| (present.north);
\draw[flow] (read.south) -- ++(0,-.23) -| (absent.north);
\node[box,text width=11.7cm] (unique) at (0,-9.15)
  {Require exactly one slot $p'$ with
   $Z(\tau_{p'})\ \lor\ A_*(\tau_{p'},Q_{\tau_{p'}})$.\\
   Test every slot, including every metadata kind.};
\draw[flow] (present.south) -- ++(0,-.3) -| (unique.north);
\draw[flow] (absent.south) -- ++(0,-.3) -| (unique.north);
\node[box,text width=11.7cm] (finish) at (0,-10.5)
  {\textbf{In both branches:} check the exact end at the decoded origin.};
\draw[flow] (unique.south) -- (finish.north);
\end{scope}
\end{tikzpicture}
\caption{The two end searches and the suffix decoder. At a genuine stencil
 cut, absence of one larger search forces all smaller searches to be absent.
 The overlapping seeds then have a common first mismatch, verified by the
 reference end check in the right branch. In the left branch, every clock
 anchor is bounded data, even if the episode end comes from a later mismatch.
 Both branches test uniqueness over the entire stencil and then check the
 exact end at the decoded origin.}
\label{decode:two-guards-figure}
\end{figure}

\FloatBarrier

\begin{lemma}[Suffix decoder]\label{decode:decoder}
Suppose the suffix test just described is paired with a genuine
stencil prefix $[s,k)$, $k=s+p$, of a first episode $e\in D$ in a
word of $D^*$.
If the suffix test succeeds, it decodes the actual slot $p$ and
consumes exactly to the actual endpoint $b$.
Conversely, on the actual suffixes of a fixed episode $e\in D$,
fewer than $\mu$ actual slots fail the decoder.  In particular,
each kind has a copy at which decoding succeeds.
\end{lemma}

\begin{proof}
\emph{The true witness survives before uniqueness.} Pair the suffix
with a genuine prefix cut $k=s+p$. Its candidates are exactly
$t+p-p'$, including $t$ at $p'=p$. Every bounded inspection for these
candidates lies inside the true first episode by the margin lemma, and
is mandatory in the proposed suffix. Thus a suffix that reaches the
uniqueness test has read the genuine values of all those inspections,
regardless of whether its proposed end is too early or too late.

If one larger search is absent, the decoder requires absence of all
smaller searches and reads every seed. The true seed is among them.
The overlap argument used above makes their periodic continuations the
same, with agreement throughout their union. Consequently the reference
exact-end check identifies the true first mismatch $q_0$: an earlier
claimed mismatch contradicts agreement before $q_0$, and a later one
would fail to detect $q_0$. The reference guard therefore forces the true
end in this branch. All tested anchors agree with
\eqref{decode:anchors}, so at $t$ the exceptional predicate is true.

If instead all larger searches are present, every candidate anchor is a
bounded return, already read inside the true episode. In particular the
true candidate uses $Q_t=\zeta(t)=q$. No assumption about the proposed
end has been made or is needed to evaluate its predicate. Because the
first episode belongs to $D$, the true candidate again satisfies
\eqref{decode:unique-witness}. Notice that this branch need not have a
bounded episode end: the true smaller search can still be absent.

\emph{Uniqueness and endpoint soundness.} In either branch the actual
slot is a witness before uniqueness is imposed. Hence a successful
uniqueness test must choose that slot $p$. The final end check is now the
end rule at the true origin and forces the actual endpoint $b$. This
last check is imposed in both branches. It cannot be omitted merely
because all larger searches are present.

\emph{Completeness at almost every slot.} Use the actual suffix for each
slot of a fixed $e\in D$. The margin lemma supplies every required
window. When one larger search is absent, implication
\eqref{decode:two-searches} makes all smaller searches absent and the
common-seed argument makes the reference guard pass. For each actual slot
$p$, the computed anchors are the restriction of \eqref{decode:anchors}
to its candidates $t+p-p'$: a present $\zeta$ supplies its return,
while an absent $\zeta$ uses the common actual mismatch $q_0$.
Thus the true witness and its final end check pass. Only an additional
witness can prevent success.

For each unsuccessful actual slot $p$, select one additional witnessing
slot $p'\ne p$. Then $(p,p')$ is counted by
\eqref{decode:whole-stencil-bound}. Different unsuccessful actual slots
have different first coordinates, so this selection is injective.
Fewer than $\mu$ slots therefore fail. Since every kind has $\mu$
copies, at least one copy of each kind succeeds.
\end{proof}

The two end-search sizes are essential to this test.  If a false
origin has absent larger search, all the smaller searches,
including the true one, are absent and the shared mismatch is
already the actual endpoint anchor.  Otherwise every candidate
clock anchor is bounded data.  Thus the decoder never needs to
look for an unverified future mismatch beyond a short true
episode merely to decide a false origin's witness.
